%-------------------------
% Currículum en LaTeX — Tobias Anuf (Español)
% Template: Jake Gutierrez's resume (MIT), based on sb2nov/resume
%
% Build (from the repo root):
%   tectonic -o . resume/resume_TobiasAnuf_ES.tex
% or upload to Overleaf and compile with pdfLaTeX.
%-------------------------

\documentclass[letterpaper,11pt]{article}

\usepackage{latexsym}
\usepackage[empty]{fullpage}
\usepackage{titlesec}
\usepackage{marvosym}
\usepackage[usenames,dvipsnames]{color}
\usepackage{verbatim}
\usepackage{enumitem}
\usepackage[hidelinks]{hyperref}
\usepackage{fancyhdr}
\usepackage[english]{babel}
\usepackage{tabularx}
\usepackage{iftex}

\ifPDFTeX
  \input{glyphtounicode}
  % Ensure that the generated PDF is machine readable / ATS parsable
  \pdfgentounicode=1
\fi

\hypersetup{
  pdftitle={Tobias Anuf - Currículum},
  pdfauthor={Tobias Anuf}
}

\pagestyle{fancy}
\fancyhf{} % clear all header and footer fields
\fancyfoot{}
\renewcommand{\headrulewidth}{0pt}
\renewcommand{\footrulewidth}{0pt}

% Adjust margins
\addtolength{\oddsidemargin}{-0.5in}
\addtolength{\evensidemargin}{-0.5in}
\addtolength{\textwidth}{1in}
\addtolength{\topmargin}{-.5in}
\addtolength{\textheight}{1.0in}

\urlstyle{same}

\raggedbottom
\raggedright
\setlength{\tabcolsep}{0in}

% Sections formatting
\titleformat{\section}{
  \vspace{-4pt}\scshape\raggedright\large
}{}{0em}{}[\color{black}\titlerule \vspace{-5pt}]

%-------------------------
% Custom commands
\newcommand{\resumeItem}[1]{
  \item\small{
    {#1 \vspace{-2pt}}
  }
}

\newcommand{\resumeSubheading}[4]{
  \vspace{-2pt}\item
    \begin{tabular*}{0.97\textwidth}[t]{l@{\extracolsep{\fill}}r}
      \textbf{#1} & #2 \\
      \textit{\small#3} & \textit{\small #4} \\
    \end{tabular*}\vspace{-7pt}
}

\newcommand{\resumeSubSubheading}[2]{
    \item
    \begin{tabular*}{0.97\textwidth}{l@{\extracolsep{\fill}}r}
      \textit{\small#1} & \textit{\small #2} \\
    \end{tabular*}\vspace{-7pt}
}

\newcommand{\resumeProjectHeading}[2]{
    \item
    \begin{tabular*}{0.97\textwidth}{l@{\extracolsep{\fill}}r}
      \small#1 & #2 \\
    \end{tabular*}\vspace{-7pt}
}

\newcommand{\resumeSubItem}[1]{\resumeItem{#1}\vspace{-4pt}}

\renewcommand\labelitemii{$\vcenter{\hbox{\tiny$\bullet$}}$}

\newcommand{\resumeSubHeadingListStart}{\begin{itemize}[leftmargin=0.15in, label={}]}
\newcommand{\resumeSubHeadingListEnd}{\end{itemize}}
\newcommand{\resumeItemListStart}{\begin{itemize}}
\newcommand{\resumeItemListEnd}{\end{itemize}\vspace{-5pt}}

%-------------------------------------------
%%%%%%  RESUME STARTS HERE  %%%%%%%%%%%%%%%%%%%%%%%%%%%%

\begin{document}

%----------ENCABEZADO----------
\begin{center}
    \textbf{\Huge \scshape Tobias Anuf} \\ \vspace{1pt}
    \small +595 (971) 862-558 $|$
    \href{mailto:tobiasanuf@gmail.com}{\underline{tobiasanuf@gmail.com}} $|$
    \href{https://www.linkedin.com/in/tobias-anuf-0003942b2}{\underline{LinkedIn}} $|$
    \href{https://github.com/tobiascca}{\underline{github.com/tobiascca}} $|$
    \href{https://tobiascca.github.io/personal-portfolio/}{\underline{tobiascca.github.io/personal-portfolio}}
\end{center}

%-----------HABILIDADES TÉCNICAS-----------
\section{Habilidades Técnicas}
 \begin{itemize}[leftmargin=0.15in, label={}]
    \small{\item{
     \textbf{Lenguajes de Programación}{: Python, SQL, Java, C/C++, C\#, Swift, JavaScript, HTML/CSS} \\
     \textbf{ML y Datos}{: scikit-learn, PyTorch, TensorFlow / TF Lite, OpenCV, MediaPipe, Transformers, Pandas, NumPy} \\
     \textbf{Datos y Herramientas}{: Amazon Redshift, MySQL, MongoDB, Git, Google Cloud Platform, Unity, Raspberry Pi} \\
     \textbf{Idiomas}{: Español (Nativo), Inglés (Proficiencia / C2)}
    }}
 \end{itemize}

%-----------EDUCACIÓN-----------
\section{Educación}
  \resumeSubHeadingListStart
    \resumeSubheading
      {Universidad Politécnica Taiwán-Paraguay}{Asunción, PY}
      {Ingeniería en Informática}{Feb. 2023 -- Presente}
      \resumeItemListStart
        \resumeItem{\textit{\textbf{Cursos relevantes}: Arquitectura de Computadoras, Sistemas Operativos, Estructuras de Datos y Algoritmos, Redes de Computadoras Avanzadas, Diseño Lógico Digital con Verilog}}
      \resumeItemListEnd
    \resumeSubheading
      {National Taiwan University of Science and Technology (NTUST)}{Taipéi, TW}
      {Intercambio de 3 semestres}{Feb. 2025 -- Jun. 2026}
      \resumeItemListStart
        \resumeItem{\textit{\textbf{Cursos relevantes}: Machine Learning y Análisis de Big Data, Visión por Computadora y Aplicaciones, Minería de Datos, Minería de Texto, Redes Sociales, Sistemas de Bases de Datos, Programación en Java}}
      \resumeItemListEnd
  \resumeSubHeadingListEnd

%-----------EXPERIENCIA-----------
\section{Experiencia}
  \resumeSubHeadingListStart
    \resumeSubheading
      {Pasante de Ciencia de Datos}{Ago. 2026 -- Oct. 2026}
      {Dato Solutions S.A. -- Investigación y Desarrollo}{Asunción, PY}
      \resumeItemListStart
        \resumeItem{Desarrollé un sistema de ML que reconstruye el RUC del importador, que los datos abiertos de aduana dejaron de publicar en 2024, recuperando ``quién importa qué'' para el producto de inteligencia comercial}
        \resumeItem{Comparé cuatro enfoques sobre más de 26M de registros; un modelo TF-IDF + kNN casi duplicó la precisión top-1 del baseline de reglas (29\% $\rightarrow$ 57\%) y alcanzó 65\% top-1 / 80\% top-5 al escalar los datos}
        \resumeItem{Definí un umbral de confianza calibrado (60\% de los despachos con 86,5\% de precisión), validé los totales contra la estadística del Banco Central del Paraguay y construí un pipeline por lotes reanudable en Python + Redshift}
        \resumeItem{Corregí errores de cálculo del asistente de consultas con LLM de la empresa, reescribí sus instrucciones y desarrollé una batería automatizada de validación (invariantes + deriva) que se ejecuta periódicamente}
      \resumeItemListEnd
  \resumeSubHeadingListEnd

%-----------PROYECTOS-----------
\section{Proyectos}
    \resumeSubHeadingListStart
      \resumeProjectHeading
          {\textbf{\href{https://github.com/tobiascca/Vision-based-Autonomous-Driving-System}{Sistema de Conducción Autónoma por Visión}} $|$ \emph{Python, OpenCV, TF Lite, Raspberry Pi}}{Jun. 2025 -- Presente}
          \resumeItemListStart
            \resumeItem{Construí un auto RC con Raspberry Pi que sigue las líneas del carril en tiempo real mediante un pipeline de detección de carriles en OpenCV y un controlador de dirección tipo PID}
            \resumeItem{Entrené un detector liviano TF Lite (Edge Impulse FOMO) de señales de pare y semáforos que anula la dirección para frenar o detenerse}
          \resumeItemListEnd
      \resumeProjectHeading
          {\textbf{\href{https://github.com/tobiascca/3D-Chess-with-HandTracking}{Ajedrez 3D Multijugador por Gestos}} $|$ \emph{Python, MediaPipe, Unity, C\#}}{Abr. 2026 -- May. 2026}
          \resumeItemListStart
            \resumeItem{Juego de ajedrez en línea controlado con gestos: MediaPipe sigue el pulgar como cursor y un pellizco mueve las piezas; Python envía el seguimiento a Unity por TCP y Mirror sincroniza a dos jugadores}
          \resumeItemListEnd
      \resumeProjectHeading
          {\textbf{\href{https://github.com/tobiascca/nightly-AI-sleep-analyzer-app}{Nightly -- Analizador de Sueño On-device}} $|$ \emph{Swift, SwiftUI, SwiftData, HealthKit}}{Jul. 2026 -- Presente}
          \resumeItemListStart
            \resumeItem{App iOS privada que graba sonidos durante la noche, transcribe el habla en el dispositivo y sincroniza con Apple Health}
          \resumeItemListEnd
      \resumeProjectHeading
          {\textbf{Aplicación Web de Chat} $|$ \emph{Python, MySQL, MongoDB}}{Jun. 2025}
          \resumeItemListStart
            \resumeItem{Chat web en tiempo real para el curso de Sistemas de Bases de Datos; diseñé el esquema SQL y comparé el almacenamiento en MySQL y MongoDB}
          \resumeItemListEnd
    \resumeSubHeadingListEnd

%-----------CURSOS Y CERTIFICACIONES-----------
\section{Cursos y Certificaciones}
  \resumeSubHeadingListStart
    \resumeProjectHeading
      {\textbf{Certificados NVIDIA DLI} $|$ \emph{\href{https://learn.nvidia.com/certificates?id=1tO0Ys3ITkGJkXM3sgBKrQ}{\underline{Fundamentos de Deep Learning}}, \href{https://learn.nvidia.com/certificates?id=g-2IIRd1RO6ZWcYuaAXLtA}{\underline{Prompts para LLMs}}, Agentes RAG}}{Oct. -- Nov. 2025}
    \resumeProjectHeading
      {\textbf{Introducción a la Ciencia de Datos}, Penguin Academy $|$ \emph{Pandas, TensorFlow, Keras, PyTorch}}{Ago. 2024}
  \resumeSubHeadingListEnd

\end{document}
